% from ega3-1-fr.tex
\oldpage[0\textsubscript{III}]{30}
qui, par définition, sont bien des sous-objets gradués du complexe (simple)
défini par $K^{\bullet\bullet}$, et font donc de ce complexe un complexe
filtré~; par ailleurs, il est clair que ces filtrations sont exhaustives et
séparées.

Il correspond à chacune de ces filtrations une suite spectrale
(\hyperref[0.11.2.2-fr]{11.2.2})~; nous désignerons par
$'E(K^{\bullet\bullet})$ et $''E(K^{\bullet\bullet})$ les suites spectrales
correspondant à $(F_{\mathrm I}^p(K^{\bullet\bullet}))$ et
$(F_{\mathrm{II}}^p(K^{\bullet\bullet}))$ respectivement, dites
\emph{suites spectrales du bicomplexe $K^{\bullet\bullet}$}, et ayant toutes
deux pour aboutissement la cohomologie
$(\mathrm H^n(K^{\bullet\bullet}))$. On montre en outre (M, XV, 6) que l'on a
\[
\label{0.11.3.2.2-fr}
  'E_2^{pq}(K^{\bullet\bullet})
  =\mathrm H_{\mathrm I}^p(\mathrm H_{\mathrm{II}}^q(K^{\bullet\bullet})),
  \qquad
  ''E_2^{pq}(K^{\bullet\bullet})
  =\mathrm H_{\mathrm{II}}^q(\mathrm H_{\mathrm I}^p(K^{\bullet\bullet})).
  \tag{11.3.2.2}
\]

Tout morphisme $u:K^{\bullet\bullet}\to K^{\prime\bullet\bullet}$ de
bicomplexes est \emph{ipso facto} compatible avec les filtrations de même
type de $K^{\bullet\bullet}$ et $K^{\prime\bullet\bullet}$, donc définit un
morphisme pour chacune des deux suites spectrales~; en outre, deux morphismes
\emph{homotopes} définissent une homotopie \emph{d'ordre $\leq1$} des
complexes (simples) filtrés correspondant, donc le \emph{même} morphisme pour
chacune des deux suites spectrales (M, XV, 6.1).
\end{env}

\begin{proposition}[11.3.3]
\label{0.11.3.3-fr}
Soit $K^{\bullet\bullet}=(K^{ij})$ un bicomplexe dans une catégorie abélienne
$\mathcal C$.
\begin{enumerate}
  \item[(i)] S'il existe $i_0$ et $j_0$ tels que $K^{ij}=0$ pour $i<i_0$ ou
    $j<j_0$ (resp. $i>i_0$ ou $j>j_0$), les deux suites spectrales
    $'E(K^{\bullet\bullet})$ et $''E(K^{\bullet\bullet})$ sont
    birégulières.
  \item[(ii)] S'il existe $i_0$ et $i_1$ tels que $K^{ij}=0$ pour $i<i_0$ ou
    $i>i_1$ (resp. s'il existe $j_0$ et $j_1$ tels que $K^{ij}=0$ pour
    $j<j_0$ ou $j>j_1$) les deux suites spectrales
    $'E(K^{\bullet\bullet})$ et $''E(K^{\bullet\bullet})$ sont
    birégulières.
\end{enumerate}
Supposons en outre que dans $\mathcal C$ les limites inductives filtrantes
existent et soient exactes. Alors~:
\begin{enumerate}
  \item[(iii)] S'il existe $i_0$ tel que $K^{ij}=0$ pour $i>i_0$ (resp. s'il
    existe $j_0$ tel que $K^{ij}=0$ pour $j<j_0$), la suite
    $'E(K^{\bullet\bullet})$ est régulière.
  \item[(iv)] S'il existe $i_0$ tel que $K^{ij}=0$ pour $i<i_0$ (resp. s'il
    existe $j_0$ tel que $K^{ij}=0$ pour $j>j_0$), la suite
    $''E(K^{\bullet\bullet})$ est régulière.
\end{enumerate}
\end{proposition}

\begin{proof}
La proposition résulte aussitôt des définitions
(\hyperref[0.11.1.3-fr]{11.1.3}) et de
(\hyperref[0.11.2.4-fr]{11.2.4}), ainsi que des observations suivantes
relatives à la filtration $F_{\mathrm I}$ (et des observations analogues qu'on
en déduit pour $F_{\mathrm{II}}$ en échangeant les rôles des deux indices
dans $K^{\bullet\bullet}$)~:
\begin{enumerate}
  \item[$1^\circ$] S'il existe $i_0$ tel que $K^{ij}=0$ pour $i>i_0$, la
    filtration $F_{\mathrm I}(K^{\bullet\bullet})$ est discrète.
  \item[$2^\circ$] S'il existe $i_0$ tel que $K^{ij}=0$ pour $i<i_0$, la
    filtration $F_{\mathrm I}(K^{\bullet\bullet})$ est co-discrète. On en
    déduit aussitôt qu'il en est de même de la filtration correspondante
    $F_{\mathrm I}(\mathrm H^n(K^{\bullet\bullet}))$ pour tout $n$~; en
    outre, la définition de $B_r^{pq}$ correspondant à la filtration
    $F_{\mathrm I}(K^{\bullet\bullet})$
    (\hyperref[0.11.2.2-fr]{11.2.2}) montre que pour tout couple $(p,q)$, la
    suite $(B_r^{pq})_{r\geq2}$ est stationnaire.
  \item[$3^\circ$] S'il existe $j_0$ tel que $K^{ij}=0$ pour $j<j_0$, on a
    \[
      F_{\mathrm I}^{p+r}(K^{\bullet\bullet})
      \cap\left(\sum_{i+j=n}K^{ij}\right)=0
    \]
    dès que $p+r+j_0>n$, donc
    $Z_r^{pq}=Z_\infty(E_2^{pq})$ pour $r>q-j_0+1$~; d'autre part,
    $\mathrm H^n(F_{\mathrm I}^p(K^{\bullet\bullet}))=0$ pour
    $p>n-j_0+1$.
  \item[$4^\circ$] S'il existe $j_0$ tel que $K^{ij}=0$ pour $j>j_0$, on a
    \[
      F_{\mathrm I}^{p-r+1}(K^{\bullet\bullet})
      \cap\left(\sum_{i+j=n}K^{ij}\right)=\sum_{i+j=n}K^{ij}.
    \]
    